\chapter{Euler-Bernoulli Beam Element}
\label{ch:beam}

\section{Euler-Bernoulli Beam Element}

The Euler-Bernoulli beam element models bending with the assumption that
plane cross-sections remain plane and perpendicular to the neutral axis
(no shear deformation).

\subsection{Degrees of Freedom}

Each node has two DOFs: transverse displacement $v$ and rotation $\theta$
(= $\dd v/\dd x$). The element has 4 DOFs total:
$\vect{u}^e = [v_1,\, \theta_1,\, v_2,\, \theta_2]\transpose$.

\subsection{Hermite Shape Functions}

Cubic Hermite polynomials ensure $C^1$ continuity (continuous displacement
and slope across elements):
\begin{align}
    H_1(x) &= 1 - 3\xi^2 + 2\xi^3 \notag \\
    H_2(x) &= L(\xi - 2\xi^2 + \xi^3) \notag \\
    H_3(x) &= 3\xi^2 - 2\xi^3 \notag \\
    H_4(x) &= L(-\xi^2 + \xi^3)
    \label{eq:hermite}
\end{align}
where $\xi = x/L \in [0,1]$.

\subsection{Beam Stiffness Matrix}

The curvature is $\kappa = \dd^2 v/\dd x^2 = \matr{B}_b\, \vect{u}^e$ and the
bending moment $M = EI\kappa$. The element stiffness is:
\begin{equation}
    \matr{k}^e_b = \int_0^L EI\, \matr{B}_b\transpose \matr{B}_b \dd x
    = \frac{EI}{L^3}
    \begin{bmatrix}
         12  &  6L  & -12  &  6L  \\
         6L  &  4L^2 & -6L  &  2L^2 \\
        -12  & -6L  &  12  & -6L  \\
         6L  &  2L^2 & -6L  &  4L^2
    \end{bmatrix}
    \label{eq:beam_stiffness}
\end{equation}

\begin{importantbox}
\begin{equation*}
    \matr{k}^e_b = \frac{EI}{L^3}
    \begin{bmatrix}
         12  &  6L  & -12  &  6L  \\
         6L  &  4L^2 & -6L  &  2L^2 \\
        -12  & -6L  &  12  & -6L  \\
         6L  &  2L^2 & -6L  &  4L^2
    \end{bmatrix}
\end{equation*}
\end{importantbox}

\subsection{Consistent Load Vector (Uniform Load $q$)}

\begin{equation}
    \vect{f}^e = \frac{qL}{12}
    \begin{bmatrix} 6 \\ L \\ 6 \\ -L \end{bmatrix}
    \label{eq:beam_load}
\end{equation}

\section{Worked Example: Simply Supported Beam}

\begin{example}[Two-element simply supported beam]
\label{ex:beam2}

A beam with $L = 2$ m, $EI = 10^4$ N$\cdot$m$^2$, simply supported at
both ends, with central point load $P = 1000$ N.

\medskip
\noindent Discretize with 2 equal beam elements ($L_e = 1$ m).
\medskip

\noindent\textbf{DOFs:} node 1: $(v_1, \theta_1)$, node 2: $(v_2, \theta_2)$,
node 3: $(v_3, \theta_3)$.

\noindent\textbf{BCs:} $v_1 = 0$, $v_3 = 0$ (pin supports).

\noindent\textbf{Free DOFs:} $[\theta_1, v_2, \theta_2, \theta_3]$.

\noindent\textbf{Analytical mid-span:} $v_2 = PL^3/(48EI) = 8.33 \times 10^{-4}$ m.

\noindent See \Cref{lst:beam2d} in the Appendix for the MATLAB implementation.
\end{example}
